%-------------------------------%
%  Author: Alessandro Sciarra   %
%    Date: 02 Jul 2025          %
%-------------------------------%

%===================================================================================================%
\subsection{DRY}
%~~~~~~~~~~~~~~~~~~~~~~~~~~~~~~~~~~~~~~~~~~~~%
\begin{frame}[fragile]{Don't Repeat Yourself $\;\to\;$ The DRY principle}
    \vspace{1mm}
    \begin{overlayarea}{\textwidth}{0.8\textheight}
        \begin{itemize}
            \item If you are doing \alert{\emph{Copy \& Paste}}, you are doing it wrong!
            \item One of the most basic and yet most disregarded principle
            \item Code duplication makes the code less maintainable and easy to break
        \end{itemize}
        \begin{varblock}{block example}[11cm]{Trivial example}<only@2-3>
            \begin{onlyenv}<2>
                \begin{lstlisting}[style=MyCpp]
                    double sum1=0;
                    for(const auto& data : dataSet1)
                        sum1 += data;
                    sum1 /= dataSet1.size();


                    double sum2=0;
                    for(const auto& data : dataSet2)
                        sum2 += data;
                    sum2 /= dataSet2.size();
                \end{lstlisting}
            \end{onlyenv}
            \begin{onlyenv}<3>
                \begin{lstlisting}[style=MyCpp]
                    double calculateAverage(const std::vector<double>& dataSet)
                    {
                        double sum=0;
                        for(const auto& data : dataSet)
                            sum += data;
                        return sum / dataSet.size();
                    }

                    calculateAverage(dataSet1);
                    calculateAverage(dataSet2);
                \end{lstlisting}
            \end{onlyenv}
        \end{varblock}
        \begin{varblock}{block example}[0.98\textwidth]{Less trivial example}<only@4-5>
            \begin{lstlisting}[style=MyCpp]
                double calculateAverage(const std::vector<double>& dataSet)
                {
                    // ...
                    throw std::runtime_error("Error in \"calculateAverage\"");
                }
            \end{lstlisting}
            \begin{uncoverenv}<5>
                \begin{lstlisting}[style=MyCpp]
                    double calculateAverage(const std::vector<double>& dataSet)
                    {
                        // ...
                        using namespace std::string_literals;
                        throw std::runtime_error("Error in \""s + __func__ + "\""s);
                    }
                \end{lstlisting}
            \end{uncoverenv}
        \end{varblock}
        \PrepareURLsymbol[PB]
        \FrameRemark<5>{\texttt{\_\_func\_\_} in C++ is an implementation-defined string and is just one among \URL*{https://stackoverflow.com/a/4384825/14967071}{other possibilities}.}
        \begin{varblock}{block example}[0.98\textwidth]{Traversing a complex structure}<6->
            \begin{onlyenv}<6>
                \begin{lstlisting}[style=MyCpp]
                    // Imagine to have to do this very often
                    for(auto& line : wiki)
                    {
                        if(line.isHeader())
                            continue;
                        for(auto& letter : line)
                        {
                            if(letter.isVowel())
                            capitalize(letter);
                        }
                    }
                    // How can you avoid duplication?
                \end{lstlisting}
            \end{onlyenv}
            \begin{onlyenv}<7>
                \begin{lstlisting}[style=MyCpp]
                    void traverseAndTransform(WikiPage& wiki,
                                              std::function<void(Letter)> transform)
                    {
                        for(auto& line : wiki){
                            if(line.isHeader())
                               continue;
                            for(auto& letter : line)
                               transform(letter);
                        }
                    }

                    traverseAndTransform(wiki, /*lambda function*/);
                \end{lstlisting}
            \end{onlyenv}
        \end{varblock}
    \end{overlayarea}
\end{frame}
%~~~~~~~~~~~~~~~~~~~~~~~~~~~~~~~~~~~~~~~~~~~~%
\subsection{KISS}
%~~~~~~~~~~~~~~~~~~~~~~~~~~~~~~~~~~~~~~~~~~~~%
\begin{frame}[fragile]{Keep It Simple, Stupid $\;\to\;$ The KISS principle}
    \begin{overlayarea}{\textwidth}{\textheight}
        \begin{varblock}{quote}[0.8\textwidth]{Avoid complexity!}[A. Einstein]<only@1>
            Everything should be done as simple as possible, but not simpler.
        \end{varblock}
        \begin{onlyenv}<1>
            \begin{itemize}
                \item If you have a solution, ask yourself if there is an easier one
                \item Complexity reduces readability and maintainability
                \item \alert{Use libraries!}
                \item Do not feel cool because you did something hard in a tricky way.\\
                \PP{Feel cool when you did the same in a trivial-to-understand way!}
                \item If you really need something complicated (do you!?), then be sure everybody understands!
            \end{itemize}
        \end{onlyenv}
        \begin{varblock}{block example}[0.86\textwidth]{Simple example}<only@2-4>
            \begin{onlyenv}<2>
                \begin{lstlisting}[style=MyCpp]
                    std::vector<int> data = {1, -3, 2, 8, 4, -9};

                    auto first = data.begin(), last = data.end();
                    while (first != last) {
                        if (*first == 4) {
                            makeReportForUser();
                            break;
                        }
                        ++first;
                    }
                \end{lstlisting}
            \end{onlyenv}
            \begin{onlyenv}<3>
                \begin{lstlisting}[style=MyCpp]
                    std::vector<int> data = {1, -3, 2, 8, 4, -9};

                    auto it = std::find(data.begin(), data.end(), 4);
                    if(it != data.end())
                        makeReportForUser();
                \end{lstlisting}
            \end{onlyenv}
            \begin{onlyenv}<4>
                \begin{lstlisting}[style=MyCpp]
                    std::vector<int> data = {1, -3, 2, 8, 4, -9};

                    //Even better
                    if(isElementInVector(data, 4))
                        makeReportForUser();
                \end{lstlisting}
            \end{onlyenv}
        \end{varblock}
        \begin{onlyenv}<4>
            \vspace{-2.5mm}
            \begin{varblock}{block example}[0.86\textwidth]{}
                \begin{lstlisting}[style=MyCpp]
                    template<typename T>
                    bool isElementInVector(const std::vector<T>& data,
                                           const T& value)
                    {
                        auto it = std::find(data.begin(), data.end(), value);
                        return it != data.end();
                    }
                \end{lstlisting}
            \end{varblock}
        \end{onlyenv}
        \begin{varblock}{block example}[0.88\textwidth]{Example from real life (Quake III Arena, 1999)}<only@5-6>
            \begin{uncoverenv}<6>
                \begin{lstlisting}[style=MyCpp]
                    float Q_rsqrt(float number)
                    {
                        long i;  float x2, y;  const float threehalfs = 1.5F;

                        x2 = number * 0.5F;  y  = number;
                        // evil floating point bit level hacking
                        i  = * (long*) &y;
                        i  = 0x5f3759df - (i >> 1);  // what the fuck?
                        y  = * (float*) &i;
                        // 1st iteration
                        y  = y * (threehalfs - (x2 * y * y));
                        // 2nd iteration, this can be removed
                        // y  = y * (threehalfs - (x2 * y * y));
                        return y;
                    }
                \end{lstlisting}
            \end{uncoverenv}
        \end{varblock}
        \begin{varblock}{block example}[0.78\textwidth]{A possible refactoring}<only@7->
            \begin{uncoverenv}<8->
                \begin{lstlisting}[style=MyCpp]
                    float Q_fastInverseSqrt(float number)
                    {
                        float approximateResult
                            = calculateApproximateInverseSqrt(number);
                        return refineResult(approximateResult, number);
                    }
                \end{lstlisting}
            \end{uncoverenv}
        \end{varblock}
        \vspace{-2.5mm}
        \PrepareURLsymbol[red]
        \begin{varblock}{block example}[0.9\textwidth]{}<only@9>
            \begin{lstlisting}[style=MyCpp]
                float calculateApproximateInverseSqrt(float number)
                {
                    //Here some very low level operations are carried out.
                    //Read @|\URL*{https://en.wikipedia.org/wiki/Fast\_inverse\_square\_root\#Algorithm}{here}|@ for a detailed explanation.
                    long floatAsInteger = * (long*) &number;
                    floatAsInteger = 0x5f3759df - (floatAsInteger >> 1);
                    return * (float*) &i;
                }
            \end{lstlisting}
        \end{varblock}
        \begin{varblock}{block example}[0.9\textwidth]{}<only@10>
            \begin{lstlisting}[style=MyCpp]
                float refineResult(float result, float initialNumber)
                {
                    //Here one iteration of the Newton's method is used.
                    //Read @|\URL*{https://en.wikipedia.org/wiki/Fast\_inverse\_square\_root\#Newton's\_method}{here}|@ for a detailed explanation.
                    float refinedResult = initialNumber * result * result;
                    return (3.0F - refinedResult) * result * 0.5F;
                }
            \end{lstlisting}
        \end{varblock}
    \end{overlayarea}
    \begin{tikzpicture}[remember picture, overlay]
        \node[visible on=<5>] at (current page.center) {
\includegraphics[width=0.5\textwidth]{Quake-III-arena}};
    \end{tikzpicture}
\end{frame}
%~~~~~~~~~~~~~~~~~~~~~~~~~~~~~~~~~~~~~~~~~~~~%
\subsection{YAGNI}
%~~~~~~~~~~~~~~~~~~~~~~~~~~~~~~~~~~~~~~~~~~~~%
\begin{frame}[fragile,t]{You Aren't Gonna Need It $\;\to\;$ The YAGNI principle}
    \vspace{-3mm}
    \begin{overlayarea}{\textwidth}{0.7\textheight}
        \begin{onlyenv}<1>
            \begin{varblock*}{quote}[0.8\textwidth]{A possible definition}[\URL*{https://ronjeffries.com/xprog/articles/practices/pracnotneed/}{Ron Jeffries}]
                %Resist that impulse, every time.
                Always implement things when you actually need them, never when you just foresee that you need them.
            \end{varblock*}
            \begin{itemize}
                \item \PP{Avoid reasonably unnecessary abstractions}\\
                      \then\textbf{Yet, it is a judgment call!}
                \item \alert{Maintaining unused code is a burden}
            \end{itemize}
            \begin{center}
                \includegraphics[height=0.25\textheight]{YAGNI}
            \end{center}
        \end{onlyenv}
        \begin{varblock}{quote}[\textwidth]{Experience makes you better}[A comment on Reddit]<only@2->
            \footnotesize
            Here's the truth with these patterns/paradigms.
            A junior dev will ignore them and code something simple that will not scale.
            A mid level will understand abstraction but over engineer a lot of things because they lack the experience to judge whether they are appropriate in the context of their work.
            Seniors will go back to simple things and postpone until it's really needed to make it more complex, because life taught them the costs of premature optimizations.
            The final level is when you understand that all of these are concepts that are good to know but none is a universal truth.
            [...]
        \end{varblock}
        \vspace{-1mm}
        \begin{onlyenv}<3->
            \begin{varblock}{example}[\textwidth]{I need to check whether a vector contains an element:}<only@3>
                \begin{lstlisting}[style=MyCpp]
                    // A (too?) specific implementation
                    bool isElementInVector(const std::vector<double>& data,
                                           double value)
                    {
                        auto it = std::find(data.begin(), data.end(), value);
                        return it != data.end();
                    }
                \end{lstlisting}
            \end{varblock}
            \begin{varblock}{example}[\textwidth]{I need to check whether a vector contains an element:}<only@4>
                \begin{lstlisting}[style=MyCpp]
                    template<typename T>
                    bool isElementInVector(const std::vector<T>& data,
                                           const T& value)
                    {
                        auto it = std::find(data.begin(), data.end(), value);
                        return it != data.end();
                    }
                \end{lstlisting}
            \end{varblock}
            \begin{varblock}{example}[\textwidth]{I need to check whether a vector contains an element:}<only@5>
                \begin{lstlisting}[style=MyCpp]
                    // Probably NOT straight away!
                    template<typename Container, typename Item>
                    bool contains(const Container& container, const Item& value)
                    {
                        // Some implementation possibly calling container
                        // member find or ::std::find if container has none
                    }
                \end{lstlisting}
            \end{varblock}
        \end{onlyenv}
    \end{overlayarea}
    \FrameRemark<1>{Ron Jeffries is a co-founder of extreme programming (XP) philosophy.}
\end{frame}
%~~~~~~~~~~~~~~~~~~~~~~~~~~~~~~~~~~~~~~~~~~~~%
\subsection{Optimisation}
%~~~~~~~~~~~~~~~~~~~~~~~~~~~~~~~~~~~~~~~~~~~~%
\begin{frame}{Do not be disappointed}{I have only a couple of slides about optimisation}
    \begin{varblock}{quote}[0.8\textwidth]{}
        \normalfont\Large\PS{\textbf{Correctness}} comes before \alert{performance}!
    \end{varblock}
    \begin{varblock*}{}[0.72\textwidth]{Steps needed seeking for speed:}
        \begin{enumerate}
            \item Write \textbf{enough} tests
            \item Profile the code
            \item Benchmark the part you want optimise
            \item Change the code \uncover<2>{{\small$\;\leftarrow\;$ \alert{\textbf{Optimisation is not just this!}}}}
            \item Make sure tests still pass
            \item \CPP|goto| \;\MakeEnumerateBox{3} or \MakeEnumerateBox{2}
        \end{enumerate}
    \end{varblock*}
\end{frame}
%~~~~~~~~~~~~~~~~~~~~~~~~~~~~~~~~~~~~~~~~~~~~%
\begin{frame}{Beware of optimisations}
    \vspace{-5mm}
    \begin{quoteblock}{Premature optimisation}[D. Knuth]
        Programmers waste enormous amounts of time thinking about, or worrying about,
        the speed of noncritical parts of their programs, and these attempts at efficiency
        actually have a strong negative impact when debugging and maintenance are considered.
        We should forget about small efficiencies, say about 97\% of the time:
        \par\smallskip
        \centerline{\alert{\textbf{Premature optimisation is the root of all evil}.}}
        \par\smallskip
        Yet we should not pass up our opportunities in that critical 3\%.
    \end{quoteblock}
    \vspace{-3mm}
    \begin{itemize}
        \item Compiler do more and more a better job!
        \item Use optimisations flags (e.g. \PP{\texttt{-O1}}, \PP{\texttt{-O3}} with \texttt{g++} and \texttt{clang})
        \item \alert{Never} sacrifice clean code in name of \textbf{claimed} better performance!
        \item Only optimise after having used a profiler and having written tests!
    \end{itemize}
\end{frame}
%~~~~~~~~~~~~~~~~~~~~~~~~~~~~~~~~~~~~~~~~~~~~%
\subsection{Boy scout philosophy}
%~~~~~~~~~~~~~~~~~~~~~~~~~~~~~~~~~~~~~~~~~~~~%
\begin{frame}{The boy scout rule}
    \vspace{0mm}
    \begin{columns}[c]
        \hfill
        \begin{column}{0.46\textwidth}
            \vspace{-\medskipamount}
            \begin{varblock}{}[\columnwidth]{}
                \color{PB}Leave the campground cleaner than you found it
            \end{varblock}
            \vspace{-1mm}
            \begin{varblock}{example}[\columnwidth]{}
                \color{PS}Always check in the code cleaner than you checked it out
            \end{varblock}
        \end{column}
        \begin{column}{0.45\textwidth}
            \includegraphics[width=0.9\textwidth]{Firecamp.jpg}
        \end{column}
    \end{columns}
    \vspace{3mm}
    \begin{itemize}
        \item To clean an existing code base takes time and requires effort
        \item Aim for tiny but continuous progress
        \item Applying the boy scout rule is a way to pay back technical debt
        \item Any IDE will help you in quickly achieving easy refactoring
    \end{itemize}
    \begin{tikzpicture}[remember picture, overlay]
        \path[to] ($(current page.north west)!0.52!(current page.west)+(9mm,0)$) edge[out=180, in=180] ++(0,-18.5mm);
    \end{tikzpicture}
\end{frame}
%===================================================================================================%